\documentclass[11pt]{article}
\usepackage[T1]{fontenc}
\usepackage{lmodern}
\usepackage[margin=1in]{geometry}
\usepackage{amsmath,amssymb,amsthm,mathtools}
\usepackage{booktabs,longtable,array,microtype}
\usepackage{xcolor}
\usepackage[hidelinks]{hyperref}
\usepackage{bookmark}
\newtheorem{theorem}{Theorem}[section]
\newtheorem{proposition}[theorem]{Proposition}
\newtheorem{lemma}[theorem]{Lemma}
\newtheorem{corollary}[theorem]{Corollary}
\theoremstyle{remark}
\newtheorem{remark}[theorem]{Remark}
\newtheorem{definition}[theorem]{Definition}
\newcommand{\R}{\mathbb R}
\newcommand{\C}{\mathbb C}
\newcommand{\N}{\mathbb N}
\newcommand{\dd}{\,\mathrm d}
\newcommand{\e}{\mathrm e}
\newcommand{\sinc}{\operatorname{sinc}}
\newcommand{\diag}{\operatorname{diag}}
\DeclareMathOperator{\Var}{Var}
\DeclareMathOperator{\Rea}{Re}
\allowdisplaybreaks[2]
\setlength{\emergencystretch}{2em}
\title{Boundary-preserving comparisons for theta Toeplitz determinants\\
\large Finite certificates, positive references, and growing-dimension limitations}
\author{Edward Baker\thanks{Working draft prepared for Edward Baker with substantial assistance from OpenAI language models. The consolidation was prepared with gpt-6-astra, recorded reasoning effort ultra; contributing analytical and certificate work used the models and efforts recorded in the supporting notes. The mathematical claims remain subject to author verification and specialist review.}}
\date{4 November 2026\quad---\quad First consolidated working draft}
\begin{document}
\maketitle
\begin{abstract}
We study comparison methods for consecutive Toeplitz determinants of the
even Taylor coefficients of the Riemann xi-function. A radial Green identity
retains the normalized endpoint term while replacing the differentiated
theta source by its positive primitive. This gives a joint determinant
residual and three finite, certified lower allowances. We then compare
three reference constructions. Reciprocal finite-pole references give
positive fixed-index margins but have an exact determinant-vanishing obstruction for a fixed pole count.
An endpoint-matched Gaussian has positive minors in every dimension, yet
its determinant scale is asymptotically incompatible with the theta
coefficients, excluding uniform relative slack in a proportional regime.
Replacing each nonreal conjugate pair in the complete canonical product
by an endpoint-matched positive pair removes that growth mismatch and
gives an explicit bounded-order sector margin. However, a family of
algebraic products has the same endpoints and leading radial zero count,
a uniform weighted coefficient defect of order $q^{-2}$, and a negative
$(q,q)$ determinant. These results identify what the specified comparisons
can and cannot transfer. They establish neither a new canonical positivity
range nor the all-order positivity required by the Riemann hypothesis.
\end{abstract}

\noindent\textbf{Draft status.}
This is a consolidation of existing research, not a new numerical campaign.
The finite comparison, reciprocal comparison and Gaussian obstruction have
the internal assessment statuses recorded in Appendix~\ref{sec:provenance}.
The endpoint-pair construction and diagonal control in
Section~\ref{sec:pair-comparison} were also accepted in the completed
E04 r7 assessment: all twelve obligations passed within scope, without
mathematical correction.
No part has independent specialist human review. No priority claim is made
for the classical determinant, residue or sector mechanisms.

\clearpage
\tableofcontents
\clearpage

\section{Introduction}
Let
\begin{equation}\label{eq:intro-C}
 C(z)=\int_0^\infty\Phi(u)\cosh(u\sqrt z)\dd u
     =\frac18\xi\!\left(\frac{1+\sqrt z}{2}\right)
     =\sum_{j\ge0}c_jz^j,
 \qquad
 D_{r,k}=\det[c_{k+j-i}]_{0\le i,j<r},
\end{equation}
where $c_j=0$ for $j<0$, $D_{0,k}=1$, and
$\xi(s)=\tfrac12s(s-1)\pi^{-s/2}\Gamma(s/2)\zeta(s)$.
The apparent square root is harmless by $\xi(s)=\xi(1-s)$.
The source and its positive primitive are
\begin{align}
 \Phi(u)&=\sum_{n\ge1}
 (2\pi^2n^4e^{9u}-3\pi n^2e^{5u})e^{-\pi n^2e^{4u}},\\
 p(u)&=e^u\sum_{n\ge1}e^{-\pi n^2e^{4u}},\qquad u\ge0.
\end{align}
Every summand of $\Phi$ is positive on this half-line. We use
$b=C(1)=1/16$, $c_0=C(0)$, and $E=C/c_0$ throughout.
The total-positivity formulation of the xi zero problem is classical;
Katkova~\cite{Katkova0505174v1} discusses both that formulation and
fixed-order eventual positivity. Restating an all-order sign condition
does not supply a sign estimate.

The question here is narrower: which explicit positive references retain
enough information to pay a complete determinant residual? Endpoints,
source truncations, matrix dimension and coefficient index matter
simultaneously. We retain them in each statement rather than infer a
global induction from finite successful comparisons.

\subsection{What the paper establishes}
The Green representation produces an exact boundary contribution and a
comparison in the undifferentiated primitive. The useful finite output is
a lower bound for a complete condensation product, not merely closeness
of isolated moments. The saved certificates give
\[
 Q_{3,1}>7\cdot10^{-23},\qquad
 Q_{3,4}>10^{-74},\qquad
 Q_{4,1}>3\cdot10^{-31},\qquad
 Q_{r,k}=D_{r+1,k}D_{r-1,k}.
\]
They demonstrate a reproducible error payment at those cases. They do
not enlarge the known positivity domain.

The structural results explain why that finite success does not yet
extend to growing order. Fixed compact Taylor payments and fixed finite
pole references meet exact obstructions. A Gaussian reference avoids
finite rank, but its explicit factorial margin is too large compared
with the complete theta determinant: the normalized target tends to zero.
Finally, a complete pair reference matches both endpoints and growth up
to a bounded factor. Its scalar coefficient residual can nevertheless be
arbitrarily small while a growing diagonal determinant is negative in an
explicit algebraic control family. The last statement concerns that
larger class, not the theta source.

\subsection{Relation to existing results and evidence}
Elementary/complete Schur duality, Toeplitz minors and their partition
conventions are standard~\cite{GarciaGarciaTierz2020Accepted}.
The bounded-order angular argument is a form of Schoenberg sector
positivity, as stated in Katkova and in Khare's author
notes~\cite{Schoenberg1955Descartes,Katkova0505174v1,KhareMath341Retrieved20261104}.
We give its specific pair multiplier and margin rather than claim a new
positivity principle. A reported fixed-shift xi result also uses the
reciprocal-pole mechanism~\cite{KenanFifteenShifts2026}; its supporting
notebooks have not been audited here. Published finite-zero information
is an imported input, not a certificate rerun by this work
\cite{PlattTrudgian200409765v1}.

The conditional contour estimate in Appendix~\ref{sec:contour} is retained
to delimit a possible large-index region. Its scalar hypotheses have
not been acquired for the canonical source in this investigation. The
literature catalog already reports a stronger cubic wedge
\cite{Michalowski260716795v1}; that report is not imported as a checked
certificate. Consequently the conservative eighth-power allowance is
not presented as improved coverage or as the preferred next task.

% Prepared for Edward Baker, 4 November 2026.
% Consolidated with substantial LLM assistance: gpt-6-astra, effort ultra.
% Manuscript fragment; requires amsmath, amssymb, amsthm and the proposition environment.
% Local evidence paths are collected by the main manuscript in sec:provenance.
\section{Green transfer, complete boundary terms, and finite comparisons}
\label{sec:green-comparison}

The following calculation retains the undifferentiated theta primitive and
its inner boundary.  It produces a quantitative finite comparison, but
positivity of the primitive alone does not sign the resulting determinants.
The comparison and its limitations consolidate RGD Notes~01, 05 and~07;
their exact paths and evidence relationships are recorded in
Appendix~\ref{sec:provenance}.  Throughout this section the physical heat
time is zero, and all negative coefficient or moment indices mean zero.

\subsection{Normalization and the complete boundary expansion}
Put
\begin{align}
 p_n(u)&=\exp(u-\pi n^2e^{4u}), & p(u)&=\sum_{n\ge1}p_n(u),\\
 \Phi(u)&=\sum_{n\ge1}\bigl(2\pi^2n^4e^{9u}-3\pi n^2e^{5u}\bigr)
                    e^{-\pi n^2e^{4u}}, &
 m_j(u)&=\frac{u^{2j}}{(2j)!}\quad(j\ge0),\\
 c_j&=\int_0^\infty\Phi(u)m_j(u)\,du, &
 d_j&=\int_0^\infty p(u)m_j(u)\,du.
\end{align}
We use ordinary coefficients, so
\begin{equation}
 C(z)=\sum_{j\ge0}c_jz^j
 =\int_0^\infty\Phi(u)\cosh(u\sqrt z)\,du
 =\frac18\xi\!\left(\frac{1+\sqrt z}{2}\right),
 \qquad C(1)=b:=\frac1{16}.
 \label{eq:green-C}
\end{equation}
The even power series makes the square-root notation single-valued.
Define
\begin{equation}
 D_{r,k}=\det[c_{k+j-i}]_{0\le i,j<r},\qquad
 D_{0,k}=1,\qquad r\ge1,\ k\ge0.
 \label{eq:green-D}
\end{equation}

For clarity, the radial origin of the transfer can be checked directly.
With \(R=e^{2u}\) and
\(A(R)=\sum_{n\ge1}e^{-\pi n^2R^2}\), one has
\(p(u)=R^{1/2}A(R)\) and
\begin{equation}
 \Phi(u)=\frac12R^{5/2}\left(\partial_R^2+\frac2R\partial_R\right)A(R),
 \qquad p''-p=8\Phi.
\end{equation}
Theta reflection gives
\(1+2A(R)=R^{-1}[1+2A(1/R)]\).  Differentiating at \(R=1\)
yields \(A'(1)=-A(1)/2-1/4\), hence \(p'(0)=-1/2\).
The complete source tail makes every fixed polynomial derivative transfer
below absolutely convergent and removes the outer boundary.

\begin{proposition}[Green transfer with the inner boundary]
For every polynomial \(f\),
\begin{equation}
 \int_0^\infty\Phi f
 =b f(0)+\frac{p(0)}8f'(0)
   +\frac18\int_0^\infty p(f''-f).
 \label{eq:green-transfer}
\end{equation}
In particular,
\begin{equation}
 c_j=b\delta_{j0}+\frac18(d_{j-1}-d_j),\qquad
 c_0=b-d_0/8.
 \label{eq:green-coefficients}
\end{equation}
\end{proposition}
\begin{proof}
Integrate \((p''-p)f/8\) twice by parts.  The inner contribution is
\([-p'(0)f(0)+p(0)f'(0)]/8\).  Now use
\(m_j''=m_{j-1}\), \(m_j'(0)=0\), and \(m_j(0)=\delta_{j0}\).
\end{proof}
The tangent inequality used below also gives
\begin{equation}
 0<d_0\le\frac{e^{-\pi}}{(4\pi-1)(1-e^{-3\pi})}<\frac14,
 \qquad \frac1{32}<c_0<b.
\end{equation}
Indeed \(n^2\ge1+3(n-1)\) controls the complete sum; the last coarse
bound follows already from \(\pi>3\), \(e^3>8\) and \(e^9>2\).

The determinant version includes every lower-dimensional boundary term.
Let
\begin{equation}
 \mathcal T_{r,k}(\boldsymbol u)
 =\frac1{r!}\sum_{\sigma\in S_r}
   \det[m_{k+j-i}(u_{\sigma(i)})]_{0\le i,j<r},
\end{equation}
and, for \(0\le q\le r\), put
\begin{equation}
 W_{r,k}^{(q)}(u_0,\ldots,u_{q-1})
 =\left[\left(\prod_{i=0}^{q-1}(\partial_{u_i}^2-1)\right)
       \mathcal T_{r,k}(\boldsymbol u)\right]_{u_q=\cdots=u_{r-1}=0}.
\end{equation}
An empty integral denotes evaluation of its scalar integrand.
\begin{proposition}[Complete Green determinant identity]
For every \(r\ge1\) and \(k\ge0\),
\begin{equation}
 D_{r,k}=\sum_{q=0}^r\binom rq b^{r-q}8^{-q}
 \int_{[0,\infty)^q}\prod_{i=0}^{q-1}p(u_i)\,
                    W_{r,k}^{(q)}(\boldsymbol u)\,d\boldsymbol u.
 \label{eq:green-boundary-determinant}
\end{equation}
Moreover, \(W_{r,k}^{(q)}\equiv0\) if \(q<\min(k,r)\).
Thus only the full-volume term remains when \(k\ge r\), while
\(D_{r,0}=c_0^r\).
\end{proposition}
\begin{proof}
Expanding the row determinant and integrating against
\(\prod_i\Phi(u_i)\) gives \(D_{r,k}\); symmetrizing does not change the
integral.  Apply \eqref{eq:green-transfer} in each coordinate.
Coordinatewise evenness removes every first-derivative trace, and symmetry
makes the \(\binom rq\) terms with \(q\) volume coordinates equal.
An endpoint row has entries \(\delta_{k+j-i,0}\).  At most
\(\max(r-k,0)\) such rows can be nonzero, proving the asserted vanishing.
The triangular coefficient matrix at \(k=0\) proves the final identity.
\end{proof}
There is no further factorial in \eqref{eq:green-boundary-determinant}:
\(\mathcal T\) was defined by an average.  The probes and primitive
moment differences can have either sign.  In particular, neither
\(p>0\) nor the pure-volume representation is a positivity theorem.

\subsection{Preserving the boundary before approximating the source}
Let
\begin{equation}
 x_j=\frac18\int_0^\infty p_1m_j,\qquad
 y_j=\frac18\int_0^\infty\sum_{n\ge2}p_nm_j,\qquad
 \widehat c_j=b\delta_{j0}+x_{j-1}-x_j.
\end{equation}
Set \(a=16\pi-1\).  The tangent bound \(e^{4u}\ge1+4u\) and
\(n^2\ge4+5(n-2)\) for \(n\ge2\) give
\begin{equation}
 0<y_j\le E_j:=
 \frac{e^{-4\pi}}{8(1-e^{-5\pi})a^{2j+1}}.
 \label{eq:green-mode-tail}
\end{equation}
Indeed, the exponential tangent bounds the \(n\)-th summand by
\(e^{-\pi n^2}e^{-(4\pi n^2-1)u}\), and the normalized moment
integral equals \((4\pi n^2-1)^{-2j-1}\); summing its geometric
majorant proves \eqref{eq:green-mode-tail}.

This reference has a precise positive-measure description, but no
all-order Toeplitz positivity is assumed for it.  Write
\(\beta_n=(4\pi n^2-1)e^{-\pi n^2}/8\) and
\(\tau=b-\beta_1=\sum_{n\ge2}\beta_n>0\).
Individual-mode Green transfer shows that \((\widehat c_j)\) is the moment
sequence of \(\Phi_1(u)\,du+\tau\delta_0\).  If
\(c_0^{(1)}=\int\Phi_1\), then
\begin{equation}
 \widehat c_0-c_0=y_0,\qquad c_0-c_0^{(1)}=\tau-y_0>0,
\end{equation}
and consequently
\begin{equation}
 \frac{|c_0-c_0^{(1)}|}{|c_0-\widehat c_0|}
 \ge a^2(1-e^{-5\pi})-1>2200.
 \label{eq:green-boundary-gain}
\end{equation}
Here \(\tau\ge\beta_2=ae^{-4\pi}/8\) and
\eqref{eq:green-mode-tail} suffice; \(\pi>3\) and \(e^{15}>1000\)
give the final strict inequality.  This is a gain for the stated
\(c_0\) approximation only, not a determinant-error or runtime gain.

For a finite acquisition cutoff \(U\ge0\), put
\begin{equation}
 a_j^{[U]}=\frac1{8(2j)!}\int_0^Uu^{2j}p_1(u)\,du,
 \qquad a_U=4\pi e^{4U}-1.
\end{equation}
The complete omitted first-mode exterior obeys
\begin{equation}
 0\le x_j-a_j^{[U]}\le B_j(U):=
 \frac{e^{U-\pi e^{4U}}}{8(2j)!}
 \sum_{h=0}^{2j}\binom{2j}{h}
                \frac{U^{2j-h}h!}{a_U^{h+1}}.
 \label{eq:green-spatial-tail}
\end{equation}
To see this, write \(u=U+v\), bound
\(p_1(U+v)\le e^{U-\pi e^{4U}}e^{-a_Uv}\), expand
\((U+v)^{2j}\), and integrate.  This does not pay quadrature uncertainty
on \([0,U]\), which must be enclosed separately.  Nor is a Green
identity applied on the truncated interval.

\subsection{The combined determinant payment}
Fix \(U=2\), write \(a_j=a_j^{[2]}\), and set
\begin{equation}
 T_j=E_j+B_j(2),\qquad
 \mathcal C_j(z)=b\delta_{j0}+a_{j-1}-a_j+z_{j-1}-z_j
 \quad(j\ge0).
\end{equation}
There is an actual vector \(z\) with \(0\le z_j\le T_j\) such that
\(c_j=\mathcal C_j(z)\).  Put \(c_j^*=\mathcal C_j(0)\), and define
\begin{align}
 Q_{r,k}&=D_{r,k}^2-D_{r,k-1}D_{r,k+1}
          =D_{r+1,k}D_{r-1,k},\\
 P_{r,k}(z)&=D_{r,k}(\mathcal C(z))^2
       -D_{r,k-1}(\mathcal C(z))D_{r,k+1}(\mathcal C(z)).
 \label{eq:green-joint-Q}
\end{align}
The product identity is Desnanot--Jacobi, with all boundary zeros retained.
After collecting the joint polynomial, write
\begin{equation}
 P_{r,k}(z)=q^*_{r,k}+\sum_{\alpha\ne0}p_\alpha z^\alpha,
 \qquad \mathcal R_{r,k}=\sum_{\alpha\ne0}|p_\alpha|T^\alpha,
 \qquad L_{r,k}=q^*_{r,k}-\mathcal R_{r,k}.
 \label{eq:green-Taylor-payment}
\end{equation}
Only moments through \(k+r\) and degrees through \(2r\) can occur.
The exact finite expansion immediately proves
\begin{equation}
 |Q_{r,k}-q^*_{r,k}|\le\mathcal R_{r,k},\qquad Q_{r,k}\ge L_{r,k}.
 \label{eq:green-absolute-payment}
\end{equation}
The neighboring products must be subtracted and collected before absolute
values are taken.  The remaining payment still discards correlations
between the actual residual moments; it is sufficient, not necessary.

A useful bound on this same payment avoids listing every translated term.
Let \(H_{n,k}(c^*)\) be the collected determinant polynomial with its
coefficients replaced by their absolute values, evaluated at positive
reference coefficients.  If each nonzero used coefficient satisfies
\(\|\mathcal C_j-c_j^*\|_T\le\epsilon c_j^*\), where
\(\|\sum h_\alpha z^\alpha\|_T=\sum|h_\alpha|T^\alpha\), then
\begin{equation}
 B_{n,k}:=\|D_{n,k}(\mathcal C)-D_{n,k}(c^*)\|_T
 \le[(1+\epsilon)^n-1]H_{n,k}(c^*).
 \label{eq:green-det-norm}
\end{equation}
This follows by expanding each homogeneous degree-\(n\) monomial and
retaining every nonconstant binomial term.  Negative-index entries vanish
identically.  Coefficient convolution proves submultiplicativity of
\(\|\cdot\|_T\).  Thus, with \(d_\pm=D_{r\pm1,k}(c^*)\),
\begin{equation}
 \mathcal R_{r,k}\le
 |d_+|B_{r-1,k}+|d_-|B_{r+1,k}+B_{r+1,k}B_{r-1,k}.
 \label{eq:green-product-payment}
\end{equation}
The last term is the full product remainder.  If \(d_\pm>0\), dividing
by \(q^*_{r,k}=d_+d_-\) gives the relative allowance
\(e_++e_-+e_+e_-\), where \(e_\pm=B_{r\pm1,k}/d_\pm\).
Reference margins and conditioning therefore remain separate inputs.

\subsection{Recorded finite certificates}
\begin{proposition}[Finite quantitative comparison]
For the complete canonical source with the conventions above, the
recorded enclosures and exact arithmetic verification establish
\begin{equation}
 Q_{3,1}>7\cdot10^{-23},\qquad
 Q_{3,4}>10^{-74},\qquad
 Q_{4,1}>3\cdot10^{-31}.
 \label{eq:green-certified-Q}
\end{equation}
For the last case the unchanged compact-reference payment moreover obeys
\begin{equation}
 \frac{\mathcal R_{4,1}}{q^*_{4,1}}<\frac1{15},\qquad
 L_{4,1}>\frac7{3000}(c_1^*)^8
          >\frac7{3000^9}>3\cdot10^{-31}.
 \label{eq:green-next-order-margin}
\end{equation}
\end{proposition}
\begin{proof}[Analytical closure and certificate inputs]
For \((r,k)=(3,1),(3,4)\), the accepted record encloses the eight
compact moments \(a_0,\ldots,a_7\), all mode tails and all spatial tails,
then evaluates \eqref{eq:green-Taylor-payment} outward through degree six.
The serialized dyadic lower allowances exceed the first two rational
floors in \eqref{eq:green-certified-Q}; the saved-data verifier also
recomputes their rational propagation.  These are the finite numerical
inputs, rather than an assumption on the desired determinant sign
(RGD Note~03, the saved E02~r1 record and Review~03;
Appendix~\ref{sec:provenance}).

For the last inequality, exact verification from the same acquisition
supplies \(c_0^*,\ldots,c_5^*>0\), \(c_1^*>1/3000\), the entry-relative
budget \(\epsilon=1/40000\), and the rational boxes
\begin{equation}
 t_2\in[.465,.466],\quad t_3\in[.135,.137],\quad
 t_4\in[.028,.029],\quad t_5\in[.004,.005],\qquad
 t_j=\frac{(c_0^*)^{j-1}c_j^*}{(c_1^*)^j}.
 \label{eq:green-rational-box}
\end{equation}
All terminating decimals here denote exact rationals.  First-row
expansion at \(k=1\) gives
\begin{equation}
 D_{n,1}=\sum_{j=1}^n(-1)^{j-1}c_0^{j-1}c_jD_{n-j,1}.
\end{equation}
Writing \(D_{n,1}(c^*)=(c_1^*)^nf_n\), it follows that
\begin{align}
 f_3&=1-2t_2+t_3,\\
 f_5&=1-4t_2+3t_2^2+(3-2t_2)t_3-2t_4+t_5.
\end{align}
Endpoint monotonicity on \eqref{eq:green-rational-box} gives
\(1/5<f_3<1/4\), \(1/80<f_5<1/40\).  The corresponding collected
absolute polynomials satisfy \(S_3<21/10\), \(S_5<21/5\).
Equation \eqref{eq:green-det-norm}, including all higher binomial terms,
therefore gives
\(B_{3,1}/(c_1^*)^3<21/100000\) and
\(B_{5,1}/(c_1^*)^5<63/100000\).  Substitution in
\eqref{eq:green-product-payment} yields
\begin{equation}
 \frac{\mathcal R_{4,1}}{(c_1^*)^8}
 <\frac1{40}\frac{21}{100000}
  +\frac14\frac{63}{100000}
  +\frac{21\cdot63}{10^{10}}
 =\frac{1628823}{10^{10}}<\frac1{6000},
\end{equation}
whereas \(q^*_{4,1}>(c_1^*)^8/400\).  Subtraction proves
\eqref{eq:green-next-order-margin}.  The second exact record verifies
these input-box memberships, the polynomial identities and every rational
inequality; it acquires no additional integral
(RGD Notes~05--06, the saved E02~r2 record and Review~04;
Appendix~\ref{sec:provenance}).
\end{proof}

The first acquisition used python-flint~0.9.0 and FLINT~3.6.0 at
384-bit precision, retaining forty segment enclosures for eight moments.
The exact saved-data verifiers trust those identified integral and
constant enclosures; they do not independently reproduce the integration
or formally verify FLINT.  The analytical comparison received
cross-model internal LLM assessment.  Assessment of the first acquisition
was by its originating assistant/model with fresh same-model contributions;
the next-order argument received cross-model checking, while its new
implementation retained author checks.  None is independent specialist
review (RGD Reviews~01, 03 and~04; Appendix~\ref{sec:provenance}).

The certified quantities are the products \(Q_{r,k}\) defined in
\eqref{eq:green-joint-Q}.  No numerical normalization by a lower-order
determinant or relative induction parameter was acquired.  We assert no
new positivity range, and these finite bounds do not supply an all-order
or all-index conclusion.

\subsection{Why the successful finite payment cannot be iterated unchanged}
For the specified compact reference and independent moment bounds,
there is an exact obstruction stronger than poor conditioning of an
auxiliary estimate.  Let
\(K_0=e^{-4\pi}/[8(1-e^{-5\pi})]\), so \(E_j=K_0/a^{2j+1}\).
In the already collected joint polynomial,
\begin{equation}
 [z_{k-1}^{2r}]P_{r,k}=1,\qquad
 \mathcal R_{r,k}\ge E_{k-1}^{2r}.
\end{equation}
Indeed, its homogeneous dependence on \(z_{k-1}=v\) gives coefficient
changes \(\mathcal C_{k-1}=-v\), \(\mathcal C_k=v\); each relevant
index-\(k\) matrix is lower bidiagonal, with determinant \(v^n\).
For \(j\ge2\), the compact support gives
\(0<c_j^*\le a_0 2^{2j-2}/(2j-2)!\).  If \(k\ge r+2\),
log-convexity of the factorial along each permutation consequently gives
\begin{equation}
 |q^*_{r,k}|\le(r+1)!(r-1)!
       \left(\frac{a_0 2^{2k-2}}{(2k-2)!}\right)^{2r}.
\end{equation}
The ratio of \(E_{k-1}\) to the parenthesized quantity is
\begin{equation}
 A_k=\frac{K_0}{aa_0}\frac{(2k-2)!}{(2a)^{2k-2}},\qquad
 \frac{A_{k+1}}{A_k}=\frac{(2k)(2k-1)}{(2a)^2}\longrightarrow\infty.
\end{equation}
Thus, for every fixed \(r\ge2\),
\begin{equation}
 L_{r,k}<0\qquad\text{for all sufficiently large }k.
 \label{eq:green-exact-payment-stop}
\end{equation}
This excludes this exact compact-reference Taylor payment, not the actual
canonical determinant or a correlated source residual.  In particular,
removing the spatial cutoff changes the error problem: Note~07 proves
\(0<(c_j-c_j^{[1]})/c_j^{[1]}\le128e^{-j/(16\log j)}\) for
\(j\ge e^{16}\).  Its reference signs and conditioning still require
separate control (RGD Note~07; Appendix~\ref{sec:provenance}).  The finite certificates above
therefore furnish no all-index induction contract.

% Consolidated for Edward Baker, 4 November 2026.
% Drafting assistant: gpt-6-astra, effort ultra.
% Source: RGD Notes 10--13; Reviews 08 and 10; source ledgers 02--03.
% This is an input fragment, not a standalone manuscript.
% Requires amsmath, amssymb, amsthm, and environments proposition, lemma,
% theorem, corollary, remark.  Labels use the prefix rgd:.
% Bibliography keys used: GarciaGarciaTierz2020Accepted,
% KenanFifteenShifts2026, Odlyzko1987AuthorCopy,
% NISTDLMFPositiveThetaPrimitives.
% The enclosing manuscript should identify Edward Baker as author (or say
% drafted for him) and acknowledge substantial LLM assistance.  The internal
% E04 assessments are not independent human specialist reviews.

\section{Rectangular duality and a finite reciprocal reference}
\label{rgd:sec-reciprocal}

We retain the complete theta source and its normalization,
\begin{align}
 C(z)&=\int_0^\infty\Phi(u)\cosh(u\sqrt z)\,du
      =\frac18\xi\!\left(\frac{1+\sqrt z}{2}\right),
                                                               \label{rgd:C}\\
 C(z)&=\sum_{m\ge0}c_mz^m,\qquad C(1)=b=\frac1{16},
       \qquad c_m=0\quad(m<0),\nonumber\\
 \Phi(u)&=\sum_{n\ge1}
 (2\pi^2n^4e^{9u}-3\pi n^2e^{5u})e^{-\pi n^2e^{4u}}.
                                                               \label{rgd:source}
\end{align}
The square root in \eqref{rgd:C} creates no branch ambiguity, since
both expressions are even in it.  In particular, $c_m>0$ for $m\ge0$.
Put
\[
 \begin{gathered}
 D_{r,k}=\det[c_{k+j-i}]_{0\le i,j<r},\qquad
 E(z)=\frac{C(z)}{c_0}=\sum_{m\ge0}e_mz^m,\\
 H(z)=\frac1{E(-z)}=\sum_{m\ge0}h_mz^m,
 \end{gathered}
\]
where both coefficient sequences are extended by zero at negative indices.
The series for $H$ is initially its Taylor series at zero.

\begin{proposition}[Exact rectangular duality]
\label{rgd:duality}
For every $r,k\ge1$,
\begin{equation}
 D_{r,k}=c_0^r\det[h_{r+j-i}]_{0\le i,j<k}.
 \label{rgd:dual-equation}
\end{equation}
For $k=0$ the left side is $c_0^r$ and the dual determinant is the
empty determinant $1$.  No restriction relating $r$ and $k$ is needed.
\end{proposition}
\begin{proof}
Let $T=[e_{j-i}]_{0\le i,j<r+k}$.  It is upper triangular with diagonal
$1$.  The formal identity $E(z)H(-z)=1$ gives
\[
 (T^{-1})_{ij}=(-1)^{j-i}h_{j-i}.
\]
Take the minor with rows $I=\{0,\ldots,r-1\}$ and columns
$J=\{k,\ldots,k+r-1\}$.  Jacobi's complementary-minor identity expresses
it as $(-1)^{\sum I+\sum J}$ times the inverse minor with rows $J^c$ and
columns $I^c$, since $\det T=1$.  The first sign is $(-1)^{rk}$.
Here $J^c=\{0,\ldots,k-1\}$ and $I^c=\{r,\ldots,r+k-1\}$;
factoring the signs $(-1)^{r+j-i}$ from this inverse minor gives a
second factor $(-1)^{rk}$.  They cancel.  Multiplication by $c_0^r$
proves the result.  All entries at negative indices remain zero.
\end{proof}

This is the rectangular form of the standard elementary/complete
Jacobi--Trudi duality; see, for example,
\cite[\S2.1]{GarciaGarciaTierz2020Accepted}.  We give the finite-matrix
proof to specify the signs and the boundary convention.  The reciprocal
pole comparison below has also been reported in prior work
\cite{KenanFifteenShifts2026}; it is presented here as a method audit,
with no claim of a new fixed-index mechanism or positivity range.\footnote{
The source record attributes the web note to kenan.works, without inferring
a personal byline.  Its primary indexed text reports all-rank coverage
for the first fifteen shifts and the same reciprocal-margin mechanism.
Direct retrieval returned HTTP 502; its supporting notebooks and reported
coverage were not audited.  That report is prior-art attribution, not a
premise of the proofs given here.}

Fix an integer $k\ge1$.  The \emph{finite pole premise} used in this
section is that $C(-z)$ has exactly $k$ zeros in $|z|\le S$, all simple
and positive, with
\begin{equation}
 0<\alpha_1<\cdots<\alpha_k<\rho<S,
 \qquad\text{and no zero on }|z|=S.
 \label{rgd:packet}
\end{equation}
This is a premise about a finite disk.  We do not assume its validity
for arbitrary $k$ or radius.  Define
\begin{equation}
 A_\ell=\frac{c_0}{\alpha_\ell C'(-\alpha_\ell)},\qquad
 H(z)=\sum_{\ell=1}^k\frac{A_\ell}{1-z/\alpha_\ell}+G(z).
 \label{rgd:partial-fractions}
\end{equation}
The complete remainder $G$ is analytic in a neighborhood of $|z|\le S$.
Since $C(0)>0$ and its real zeros in this interval are simple,
successive crossings show
\begin{equation}
 \operatorname{sgn}A_\ell=(-1)^{\ell-1}.
 \label{rgd:residue-sign}
\end{equation}
The actual residue of $H$ at $\alpha_\ell$ is $-A_\ell\alpha_\ell$.

\begin{proposition}[Positive finite-pole margin and complete residual]
\label{rgd:pole-margin}
Assume \eqref{rgd:packet}.  Let $r\ge\max\{1,k-1\}$ and put
\[
 B_{ij}=\sum_{\ell=1}^k A_\ell\alpha_\ell^{-(r+j-i)},
 \quad 0\le i,j<k,\qquad
 \Delta(\alpha)=\prod_{i<j}(\alpha_j-\alpha_i).
\]
Then
\begin{equation}
 d^*_{r,k}:=\det B
 =\Delta(\alpha)^2\prod_{\ell=1}^k
       |A_\ell|\alpha_\ell^{-r-k+1}>0.
 \label{rgd:pole-determinant}
\end{equation}
Choose any $M\ge\max_{|z|=\rho}|G(z)|$, and set
\begin{gather}
 \beta_\ell=\frac{\alpha_\ell}{\rho},\qquad
 a_{\min}=\min_\ell|A_\ell|,\qquad
 u_\ell=\frac{\prod_{j\ne\ell}(1+\beta_j)}
                  {\prod_{j\ne\ell}|\beta_\ell-\beta_j|},
 \quad u_*=\max_\ell u_\ell,                              \label{rgd:pole-constants}\\
 \delta_r=\frac{k^2Mu_*^2}{a_{\min}}\beta_k^r.
 \label{rgd:pole-delta}
\end{gather}
If $\delta_r<1$, then
\begin{equation}
 D_{r,k}\ge c_0^r d^*_{r,k}(1-\delta_r)^k>0.
 \label{rgd:pole-floor}
\end{equation}
In particular, $\delta_r\le\tfrac12$ gives the explicit floor
$2^{-k}c_0^r d^*_{r,k}$.
\end{proposition}
\begin{proof}
The index restriction makes every $r+j-i$ nonnegative.  Factor $B$ into
the Vandermonde matrices $[\alpha_\ell^i]_{i,\ell}$ and
$[\alpha_\ell^{-j}]_{\ell,j}$, separated by
$\operatorname{diag}(A_\ell\alpha_\ell^{-r})$.
The inverse-node Vandermonde has determinant
\[
 (-1)^{k(k-1)/2}\Delta(\alpha)
       \prod_{\ell=1}^k\alpha_\ell^{-(k-1)}.
\]
Its sign cancels the product of the signs in
\eqref{rgd:residue-sign}, proving \eqref{rgd:pole-determinant}.

Write $G(z)=\sum_{m\ge0}g_mz^m$ and
$R_{ij}=g_{r+j-i}$.  Cauchy's coefficient estimate gives
$|g_m|\le M\rho^{-m}$.  Scale each matrix entry $(i,j)$ by
$\rho^{r+j-i}$, and denote the resulting matrices by
$\widetilde B$ and $\widetilde R$.  Each determinant acquires the
positive multiplier $\rho^{rk}$, and
$\|\widetilde R\|_\infty\le kM$.  Moreover,
\[
 \widetilde B
 =U\operatorname{diag}(A_\ell\beta_\ell^{-r})V,
 \qquad U_{i\ell}=\beta_\ell^i,\quad V_{\ell j}=\beta_\ell^{-j}.
\]
The rows of $U^{-1}$ contain the coefficients of the Lagrange
polynomials at the positive nodes $\beta_\ell$.  Their absolute row
sums are exactly $u_\ell$, so $\|U^{-1}\|_\infty=u_*$.
If $J$ reverses the columns, then
\[
 V=\operatorname{diag}(\beta_\ell^{-(k-1)})U^{\mathsf T}J,
 \qquad
 \|V^{-1}\|_\infty
 \le\sum_\ell\beta_\ell^{k-1}u_\ell\le ku_*.
\]
Consequently, for the complete preconditioned remainder,
\[
 \|\widetilde B^{-1}\widetilde R\|_\infty
 \le ku_*\frac{\beta_k^r}{a_{\min}}u_*kM=\delta_r.
\]
When $\delta_r<1$, the real matrices
$I+s\widetilde B^{-1}\widetilde R$ are nonsingular for $0\le s\le1$,
so their determinants remain positive.  Each eigenvalue of
$I+\widetilde B^{-1}\widetilde R$ has modulus at least $1-\delta_r$.
It follows that its determinant is at least $(1-\delta_r)^k$.
Use \eqref{rgd:dual-equation} and the positive determinant scaling to
obtain \eqref{rgd:pole-floor}.  No Hermitianity is assumed.
\end{proof}

For $M>0$, an explicit starting index is
\begin{equation}
 r_0=\max\left\{1,k-1,
 \left\lceil\frac{\log(2k^2Mu_*^2/a_{\min})}
                        {\log(\rho/\alpha_k)}\right\rceil\right\}.
 \label{rgd:pole-start}
\end{equation}
For $M=0$, take $r_0=\max\{1,k-1\}$.  The chosen majorant satisfies
$\delta_{r+1}=\beta_k\delta_r$, and
\[
 \frac{c_0^{r+1}d^*_{r+1,k}}{c_0^rd^*_{r,k}}
   =\frac{c_0}{\prod_{\ell=1}^k\alpha_\ell}.
\]
The first equality concerns the majorant, not an asserted equality for
the actual error norm.  The result is a conditional fixed-$k$ estimate;
it does not supply the finite pole premise as $k$ grows.

\subsection{Why a fixed number of poles cannot supply a growing reference}

\begin{proposition}[Fixed-pole obstruction, including negative indices]
\label{rgd:finite-poles}
Let
\[
 H_N(z)=\sum_{\ell=1}^N\frac{A_\ell}{1-z/\alpha_\ell}
       =\frac{P_{N-1}(z)}{Q_N(z)},\qquad
 Q_N(z)=1+q_1z+\cdots+q_Nz^N,
\]
where $\deg P_{N-1}\le N-1$.  Let $h_m^{[N]}$ be its Taylor
coefficients for $m\ge0$, extended by zero for $m<0$.  Then
\begin{equation}
 \det[h_{r+j-i}^{[N]}]_{0\le i,j<k}=0
       \qquad(r\ge N,\ k>N).
 \label{rgd:fixed-pole-zero}
\end{equation}
This includes $k>r+1$.
\end{proposition}
\begin{proof}
The identity $Q_NH_N=P_{N-1}$ gives
\[
 h_m^{[N]}+q_1h_{m-1}^{[N]}+\cdots+q_Nh_{m-N}^{[N]}=0
       \qquad(m\ge N).
\]
In row $i$ of the last column, $m=r+k-1-i\ge r\ge N$.
The preceding $N$ columns exist because $k>N$.  Thus the last column
plus the stated linear combination of those columns is identically
zero, with the zero negative-index convention unchanged.
\end{proof}

When $r\ge k-1$ one may also factor the whole matrix through an
$N$-dimensional exponential space and conclude rank at most $N$.
That rank argument must not be used for $k>r+1$: the exponential
formula does not extend to the artificially zero negative coefficients.
The recurrence proof establishes \eqref{rgd:fixed-pole-zero} there.
The proposition concerns the reference; residual terms can still make
the complete determinant nonzero.

\subsection{A distinct limitation of the triangle--Harnack allowance}
\label{rgd:harnack-section}

There is also an explicit, but deliberately lossy, way to pay $M$ in
Proposition~\ref{rgd:pole-margin}.  Under \eqref{rgd:packet}, put
\[
 F_0(z)=\frac{C(-z)}{c_0\prod_{\ell=1}^k(1-z/\alpha_\ell)},
 \quad
 L=\frac{C(S)}{c_0\prod_{\ell=1}^k(S/\alpha_\ell-1)},
 \quad U_0\ge\log L.
\]
The function $F_0$ is zero-free on the closed disk, $F_0(0)=1$, and
$|F_0|\le L$ on its boundary.  Thus $L\ge1$.  The analytic logarithm
with $\log F_0(0)=0$, followed by Harnack's inequality for
$U_0-\operatorname{Re}\log F_0$, gives
\begin{equation}
 -\frac{2U_0v}{S-v}\le\operatorname{Re}\log F_0(z)
       \le\frac{2U_0v}{S+v}\qquad(|z|=v<S).
 \label{rgd:harnack-log}
\end{equation}
Let
\begin{gather*}
 p_\ell=\prod_{j\ne\ell}|1-\alpha_\ell/\alpha_j|,
 \quad a^-=\min_\ell A_\ell^-,\\
 A_\ell^-=\frac{e^{-2U_0\alpha_\ell/(S+\alpha_\ell)}}{p_\ell},
 \quad
 A_\ell^+=\frac{e^{2U_0\alpha_\ell/(S-\alpha_\ell)}}{p_\ell},
 \\
 \widehat M=\frac{e^{2U_0\rho/(S-\rho)}}
                      {\prod_\ell(\rho/\alpha_\ell-1)}
       +\sum_\ell\frac{A_\ell^+}{\rho/\alpha_\ell-1},
 \qquad
 \widehat\delta_r=\frac{k^2\widehat M u_*^2}{a^-}\beta_k^r.
\end{gather*}
Indeed, the residue formula and \eqref{rgd:harnack-log} give
$A_\ell^-\le|A_\ell|\le A_\ell^+$; the triangle inequality on
$G=H-\sum A_\ell/(1-z/\alpha_\ell)$ gives $\max_{|z|=\rho}|G(z)|\le\widehat M$, so one may choose
$M=\widehat M$.

\begin{proposition}[Lower bound on this particular allowance]
\label{rgd:harnack-stop}
The specified allowance satisfies
\begin{equation}
 \widehat\delta_r\ge
       4^{k-1}\frac{\beta_k^{r+1}}{1-\beta_k}.
 \label{rgd:harnack-lower}
\end{equation}
Consequently it diverges along any sequence of valid packets with
$k\to\infty$, $r\ge k-1$, $k/r$ bounded below by a positive constant,
and $\beta_k\to1$.
\end{proposition}
\begin{proof}
Interpolate the monic polynomial $(x-\tfrac12)^{k-1}$ at the
$\beta_\ell\in(0,1)$.  Comparing leading coefficients and using
$|\beta_\ell-\tfrac12|\le\tfrac12$ gives
\[
 1\le2^{1-k}\sum_\ell
       \frac1{\prod_{j\ne\ell}|\beta_\ell-\beta_j|}
    \le2^{1-k}ku_*.
\]
Also, the last residue term in $\widehat M$ implies
$\widehat M/a^-\ge1/(\rho/\alpha_k-1)=\beta_k/(1-\beta_k)$.
These inequalities prove \eqref{rgd:harnack-lower}.
For the final assertion, divide its logarithm by $k$.
The contribution $(r+1)\log\beta_k/k$ tends to zero, while
$(k-1)\log4/k$ tends to $\log4$ and
$-\log(1-\beta_k)/k\ge0$.
\end{proof}

For the canonical function, the ordinary zero-counting asymptotic
forces the last condition along any unbounded sequence of valid
packets.  To state this precisely, enumerate the moduli $R_j$ of the
zeros of $C(-z)$ with multiplicity.  A nontrivial zero
$\sigma+i\gamma$ of $\xi$, with $\gamma>0$, maps to
$\alpha=-(2\sigma-1+2i\gamma)^2$, so
$|\alpha|=4\gamma^2+(2\sigma-1)^2$.  The standard zero count
\cite[\S2]{Odlyzko1987AuthorCopy} and $0<\sigma<1$ give
\[
 n(R):=\#\{j:R_j\le R\}\sim\frac{\sqrt R\log R}{8\pi}.
\]
If $R_{j+1}/R_j$ failed to tend to one, two fixed relative interior
radii in a recurring gap would have equal counts but asymptotically
different values under this formula.  Thus $R_{j+1}/R_j\to1$.
A valid packet has $R_k=\alpha_k<\rho<S<R_{k+1}$, whence
$\beta_k\to1$.  The existence of arbitrarily large valid packets is
not assumed; it would in particular impose global reality and
simplicity on the zeros.

This is a failure of the defined allowance, not a lower bound for the
actual remainder.  For example,
$E(z)=\prod_{\ell=1}^k(1+z/\alpha_\ell)$ has exact reciprocal partial
fractions and $G=0$, whereas $\widehat M$ above is positive and can be
large.  Neither Proposition~\ref{rgd:finite-poles} nor
Proposition~\ref{rgd:harnack-stop} excludes every finite-pole argument
or signs a growing-index canonical determinant.

\section{A positive Gaussian reference and its growth obstruction}
\label{rgd:sec-gaussian}

We now consider a reference with a nonzero determinant at every
$r,k$.  Unlike the fixed-pole reference, it needs no finite-zero input.
The primitive for the complete source is
\[
 \begin{gathered}
 p(u)=\sum_{n\ge1}e^{u-\pi n^2e^{4u}},\qquad
 p''-p=8\Phi,\quad p'(0)=-\tfrac12,\\
 c_0=b-\frac{d_0}{8},\qquad d_0=\int_0^\infty p(u)\,du.
 \end{gathered}
\]
These are the canonical Green normalization identities; the derivative
at zero follows from the theta functional equation.  The following
bounds are elementary consequences of the complete sums.

\begin{lemma}[Complete-source envelope]
\label{rgd:envelope}
For $u\ge0$,
\begin{equation}
 0<\Phi(u)\le4\pi^2e^{9u-\pi e^{4u}},\qquad
 \frac1{32}<c_0<\frac1{16}.
 \label{rgd:envelope-equation}
\end{equation}
Writing $a_R=(9+\sqrt R)/4$, for every $R>0$ one has
\begin{equation}
 E(R)<\mathcal F(R):=32\pi^{2-a_R}\Gamma(a_R),
 \qquad \log\mathcal F(R)=O(\sqrt R\log R)\quad(R\ge2)
 \label{rgd:entire-growth}
\end{equation}
as an upper growth estimate.
\end{lemma}
\begin{proof}
Write $v=\pi e^{4u}$, and let $\Phi_n$ be the $n$th summand of
\eqref{rgd:source}.  Every summand is positive, and for $n\ge2$,
\[
 \frac{\Phi_n}{\Phi_1}
 =\frac{n^2(2n^2v-3)}{2v-3}e^{-(n^2-1)v}
 \le2n^4e^{-(n^2-1)v}.
\]
The elementary bounds $n^4\le16^{n-1}$ and
$n^2-4\ge5(n-2)$ show
\[
 \sum_{n\ge2}\frac{\Phi_n}{\Phi_1}
 \le\frac{32e^{-3v}}{1-16e^{-5v}}\le64e^{-3v}<1
       \qquad(v\ge\pi).
\]
Since $\Phi_1\le2\pi^2e^{9u-\pi e^{4u}}$, this proves the envelope.
Also $e^{4u}\ge1+4u$ gives
\[
 0<d_0\le\sum_{n\ge1}\frac{e^{-\pi n^2}}{4\pi n^2-1}
 \le\frac{e^{-\pi}}{(4\pi-1)(1-e^{-3\pi})}<\frac14.
\]
The Green normalization gives the asserted bounds on $c_0$.
For real $R>0$, bound the hyperbolic cosine by $e^{u\sqrt R}$ and
substitute $v=\pi e^{4u}$ in the source integral to obtain
\[
 E(R)\le\frac{\pi^{2-a_R}}{c_0}
       \int_\pi^\infty v^{a_R-1}e^{-v}\,dv
       <32\pi^{2-a_R}\Gamma(a_R).
\]
For completeness, if $N=\lceil a_R\rceil$, then
$v^{a_R-1}\le1+v^N$ for all $v>0$, and hence
$\Gamma(a_R)\le1+N!$.  The estimate $\log N!\le N\log N$ proves
the upper growth assertion without an asymptotic gamma formula.
\end{proof}

Define
\begin{equation}
 t=\log(b/c_0)>0,\qquad C_G(z)=c_0e^{tz},\qquad
 \Phi_G(u)=\frac{c_0}{\sqrt{\pi t}}e^{-u^2/(4t)}.
 \label{rgd:gaussian}
\end{equation}
The elementary Gaussian integral gives the cosh transform $C_G$,
so $C_G(0)=c_0$ and $C_G(1)=b$.  The boundary-compatible primitive
\[
 p_G(u)=8\int_u^\infty\sinh(v-u)\Phi_G(v)\,dv
\]
satisfies $p_G''-p_G=8\Phi_G$ and $p_G'(0)=-8b=-1/2$.
Tonelli's theorem gives
$\int_0^\infty p_G=8\int_0^\infty(\cosh v-1)\Phi_G(v)\,dv
=8(b-c_0)=d_0$.  Thus this comparison retains $C(0)$, $C(1)$, the derivative
$p'(0)$, and the primitive integral.  No equality of $p_G(0)$
and $p(0)$ is asserted.
The complete signed difference
\[
 \Delta E(z)=\frac{C(z)-C_G(z)}{c_0},\qquad
 \Delta E(0)=\Delta E(1)=0
\]
has coefficients
\begin{equation}
 \varepsilon_m=e_m-\frac{t^m}{m!}
 =\frac1{c_0(2m)!}\int_0^\infty
       (\Phi-\Phi_G)(u)u^{2m}\,du\quad(m\ge0).
 \label{rgd:complete-difference}
\end{equation}
They are set to zero for $m<0$.  All source modes and tails remain in
this difference; no bound on its two summands separately is used below.

\begin{proposition}[Gaussian reference margin at every size]
\label{rgd:gaussian-margin}
Let $g_m=t^m/m!$ for $m\ge0$ and $g_m=0$ for $m<0$, and define
$r$-square real matrices
\[
 A=[e_{k+j-i}]_{0\le i,j<r},\qquad
 B=[g_{k+j-i}]_{0\le i,j<r},\qquad K=B^{-1}(A-B).
\]
For every $r,k\ge1$,
\begin{equation}
 \begin{gathered}
 m_{r,k}:=\det B
   =t^{rk}\prod_{j=0}^{r-1}\frac{j!}{(k+j)!}
   \ge\left(\frac{t}{r+k}\right)^{rk}>0,\\
 \frac{D_{r,k}}{c_0^rm_{r,k}}=\det(I+K).
 \end{gathered}
 \label{rgd:gaussian-margin-equation}
\end{equation}
\end{proposition}
\begin{proof}
Factor $t^{k+j}/(k+j)!$ from column $j$ and $t^{-i}$ from row $i$.
The remaining entry is the falling factorial $(k+j)_i$, a monic
polynomial of degree $i$ evaluated at $k+j$.  Its evaluation determinant
is the Vandermonde at $k,k+1,\ldots,k+r-1$, namely
$\prod_{j=0}^{r-1}j!$.  If $i>k+j$, the falling factorial contains a
zero and exactly reproduces the negative-index convention.  The total
power of $t$ is $rk$.  Each quotient $(k+j)!/j!$ consists of $k$
factors at most $r+k$, giving the lower bound.  The final equality
is the exact determinant identity, using $D_{r,k}=c_0^r\det A$.
\end{proof}

The explicit determinant is standard algebra (equivalently a
rectangular Schur specialization); no new determinant identity is
asserted.  It removes the fixed-pole determinant-vanishing obstruction of
Proposition~\ref{rgd:finite-poles}.  Its scale nevertheless prevents a
uniform positive relative slack for the actual complete residual.

\begin{theorem}[Growth obstruction to uniform Gaussian slack]
\label{rgd:gaussian-stop}
For every $r,k\ge1$ and $R>0$,
\begin{align}
 \left|\frac{D_{r,k}}{c_0^rm_{r,k}}\right|
 &\le r^{r/2}E(R)^r\left(\frac{r+k}{tR}\right)^{rk}
 \le\overline U_{r,k}(R),                              \label{rgd:ratio-bound}\\
 \overline U_{r,k}(R)
 &:=r^{r/2}\mathcal F(R)^r
                         \left(\frac{r+k}{tR}\right)^{rk}.
 \nonumber
\end{align}
Fix $M>0$ and set $R_k=k^2/(\log k)^2$ for $k\ge3$.  Then, uniformly
for $1\le r\le Mk$ as $k\to\infty$,
\begin{equation}
 \frac1r\log\overline U_{r,k}(R_k)
 \le-k\log k+2k\log\log k+C_Mk,
 \qquad \overline U_{r,k}(R_k)^{1/r}\longrightarrow0.
 \label{rgd:ratio-decay}
\end{equation}
For every induced matrix norm of the exact residual $K$,
\begin{equation}
 \|K\|\ge1-\overline U_{r,k}(R_k)^{1/r}.
 \label{rgd:norm-stop}
\end{equation}
Thus no bound $\|K\|\le\delta_*<1$ can hold uniformly on this domain.
\end{theorem}
\begin{proof}
Positivity of the coefficients gives $e_mR^m\le E(R)$ for $m\ge0$.
Multiply entry $(i,j)$ of $A$ by $R^{k+j-i}$.  This is an invertible
row and column scaling, with determinant multiplier $R^{rk}$.  All
scaled entries have modulus at most $E(R)$; negative-index entries
remain zero.  Hadamard's inequality gives
\[
 |\det A|\le r^{r/2}E(R)^rR^{-rk}.
\]
Divide by the lower bound for $m_{r,k}$ in
\eqref{rgd:gaussian-margin-equation} and use
Lemma~\ref{rgd:envelope} to get \eqref{rgd:ratio-bound}.
At $R_k$, the upper growth bound gives
$\log\mathcal F(R_k)\le Ck$ for all sufficiently large $k$.
Since $r+k\le(M+1)k$,
\[
 \begin{aligned}
 \frac1r\log\overline U_{r,k}(R_k)
 &=\frac12\log r+\log\mathcal F(R_k)\\
 &\quad+k\bigl(\log(r+k)-\log t-\log R_k\bigr)\\
 &\le-k\log k+2k\log\log k+C_Mk.
 \end{aligned}
\]
All constants are independent of $r$ in the stated domain.  This
proves \eqref{rgd:ratio-decay}, and also uniform convergence of
$\overline U_{r,k}(R_k)$ itself to zero because $r\ge1$.
If $\|K\|<1$, every eigenvalue of $I+K$ has modulus at least
$1-\|K\|$.  The determinant identity consequently gives
\[
 (1-\|K\|)^r\le|\det(I+K)|
      \le\overline U_{r,k}(R_k),
\]
which proves \eqref{rgd:norm-stop}.  If $\|K\|\ge1$, that inequality
is immediate.
\end{proof}

The theorem includes fixed $r$, $r=k$, and any fixed proportional
ratio with $k/r\ge1/M$.  It does not cover fixed $k$ with $r\to\infty$.
It is a lower bound $\|K\|\ge1-o(1)$, not a claim that $\|K\|$ tends
to one.  If $A,B$ undergo common invertible scalings $A\mapsto LAR$
and $B\mapsto LBR$, then $K\mapsto R^{-1}KR$.  Norm values may change,
but the determinant-based obstruction remains valid for every induced
norm of the transformed residual.

There is a corresponding exact scalar obstruction:
\begin{equation}
 \frac{\det A-m_{r,k}}{m_{r,k}}\longrightarrow-1,
 \qquad
 \frac{|\det A-m_{r,k}|}{m_{r,k}}\longrightarrow1
 \label{rgd:scalar-stop}
\end{equation}
uniformly on the same domain.  It excludes a fixed positive relative
slack, not every tolerance approaching one.  Once $|\det A|<m_{r,k}$,
the bare inequality $|\det A-m_{r,k}|<m_{r,k}$ is equivalent to
$\det A>0$.  A proof using such a shrinking tolerance must therefore
supply an independently established signed estimate; merely imposing
the inequality restates the target.

The scale mismatch can also be made explicit.  If a positive moving
slope $t_{r,k}$ in the exponential family permitted a fixed norm bound
$\delta<1$, the same argument would require
\begin{equation}
 t_{r,k}\le
 \frac{r+k}{R_k(1-\delta)^{1/k}}
       r^{1/(2k)}E(R_k)^{1/k}
 =O_{M,\delta}\!\left(\frac{\log^2 k}{k}\right).
 \label{rgd:moving-slope}
\end{equation}
This is a necessary condition only.  Matching both endpoint values
fixes $t=\log(b/c_0)>0$, so it cannot meet this vanishing scale in the
same single-exponential family.

\subsection{The endpoint-preserving quadratic adjustment}

\begin{proposition}[Failure of universal positivity for a nonzero correction]
\label{rgd:quadratic}
For $q\in\mathbb R$, put
\[
 E_2(z)=\exp\bigl((t+q/2)z-qz^2/2\bigr),\qquad C_2=c_0E_2.
\]
Then $C_2(0)=c_0$ and $C_2(1)=b$.  For every $q\ne0$, some
rectangular coefficient determinant of $C_2$ is negative.
\end{proposition}
\begin{proof}
If $q>0$, then $E_2(x)\to0$ as $x\to+\infty$.  An entire series
with constant coefficient one and all coefficients nonnegative is
at least one on that axis.  Thus some coefficient $e_j^{[2]}$ is
negative, giving $D_{1,j}^{[2]}=c_0e_j^{[2]}<0$.
If $q<0$, the dual series
\[
 H_2(z)=\frac1{E_2(-z)}
       =\exp\bigl((t+q/2)z+qz^2/2\bigr)
\]
tends to zero on the positive axis, so some $h_j^{[2]}<0$.
Proposition~\ref{rgd:duality}, which is purely algebraic, gives
$D_{j,1}^{[2]}=c_0^jh_j^{[2]}<0$.
\end{proof}

Matching the first coefficient as well would select $q=2(e_1-t)$;
no sign or nonzero assertion about that canonical number is needed.
For $q=0$, Theorem~\ref{rgd:gaussian-stop} applies.  The nonzero
quadratic result concerns positivity at \emph{all} indices and uses
edge-family witnesses; it does not exclude every restricted
proportional-index subdomain or other reference families.

\subsection{A positive control and the extent of the conclusions}

The Gaussian obstruction is compatible with strict positivity of every
target rectangle.  Indeed, for $\ell>0$ take
\[
 C_*(z)=\frac{\cosh(\ell\sqrt z)}{16\cosh\ell},\qquad
 E_*(z)=\cosh(\ell\sqrt z).
\]
The classical product \cite[Eq.~4.36.2]{NISTDLMFPositiveThetaPrimitives}
gives $E_*(z)=\prod_{j\ge1}(1+x_jz)$, where
$x_j=4\ell^2/((2j-1)^2\pi^2)>0$ and $\sum_jx_j<\infty$.
For a finite product, its upper Toeplitz matrix in dimension $r+k$ is
the product of upper bidiagonal matrices $I+x_jS$, with $S$ the upper
shift.  Each factor has nonnegative minors.  In the Cauchy--Binet
expansion for rows $0,\ldots,r-1$ and columns $k,\ldots,k+r-1$,
move the full consecutive block by one at each of the first $k$
factors and leave it fixed thereafter.  This supplies the positive
term $(x_1\cdots x_k)^r$.  Local convergence of the product preserves
that lower bound in every fixed rectangle.  Thus all these target
determinants are strictly positive, while the same proof of
Theorem~\ref{rgd:gaussian-stop}, using
$\log E_*(R)\le\ell\sqrt R$, gives the Gaussian fixed-slack obstruction.
This is an atomic source control, not the canonical smooth theta source
or a certified positive-time collision.

The conclusions here are deliberately separate.  A fixed finite pole
reference eventually has a zero determinant; the specified
triangle--Harnack allowance can diverge despite a zero actual error;
and the endpoint-matched Gaussian reference has a positive margin at
every size but cannot support uniform positive relative slack on the
stated growing-index domain.  None of these facts makes a canonical
determinant negative, supplies a new theta positivity range, or proves
an all-index induction.  Continuation would require an appropriately
scaled reference with a paid complete residual, or an independently
controlled signed comparison.  No such estimate is asserted here.

\paragraph{Assessment provenance.}
The reciprocal results were accepted within their recorded scope in
internal assessment RGD-E04 r5, and the Gaussian and quadratic results
were accepted within their recorded scope in RGD-E04 r6, without a
mathematical correction.  Those assessments used a different LLM model
with shared source, algebra, and envelope premises; they were not
independent human specialist reviews.  The present consolidation adds
no numerical evidence, certificate replay, new positivity coverage, or
priority claim.  The direct algebraic proofs above are the mathematical
support for the statements, not the assessment labels.

% Prepared for Edward Baker, 4 November 2026.
% Substantial LLM assistance: gpt-6-astra, actual effort ultra.
% Optional manuscript section, based on frozen RGD Note14.
% Accepted within scope by separate internal RGD-E04 r7 assessment; no specialist review.

\section{Endpoint-preserving pair comparison and a diagonal obstruction}
\label{sec:pair-comparison}

For a generating function $F(z)=\sum_{j\ge0}f_jz^j$, with $f_0=1$,
write
\[
 \Delta_{r,k}(F)=\det[f_{k+j-i}]_{0\le i,j<r},\qquad f_j=0\quad(j<0).
\]
The original normalization is $D_{r,k}=c_0^r\Delta_{r,k}(E)$.
The comparison below preserves $E(0)$ and $E(1)$ and has the correct
entire-function growth. The subsequent obstruction concerns only the
listed algebraic data; it is not a negative determinant for the theta
source.

\begin{proposition}[Endpoint-preserving pair comparison]
\label{prop:endpoint-pair}
Suppose
\[
 E(z)=\prod_i(1+a_i z)\prod_\ell
       (1+2s_\ell\cos\theta_\ell\,z+s_\ell^2z^2),
\]
where $a_i,s_\ell>0$, $0<\theta_\ell<\pi/2$, the complete alphabet is
infinite, $\sum_i a_i+\sum_\ell s_\ell<\infty$, and
$\sum_\ell\tan^2\theta_\ell<\infty$.
Set
\[
 y_\ell=\sqrt{1+2s_\ell\cos\theta_\ell+s_\ell^2}-1,
 \quad d_\ell=s_\ell^2-y_\ell^2,
 \quad
 R(z)=\prod_i(1+a_i z)\prod_\ell(1+y_\ell z)^2.
\]
Assume also $y_\ell\le2$. Then $R(0)=E(0)=1$, $R(1)=E(1)$, and,
with
\[
 V=\prod_\ell\frac{s_\ell^2}{y_\ell^2}<\infty,
\]
one has
\[
 1\le\frac{E(t)}{R(t)}\le V\quad(t\ge1),
 \qquad \lim_{t\to\infty}\frac{E(t)}{R(t)}=V.
\]
If $E-R=\sum_j\varepsilon_jz^j$, then
\begin{equation}
 \sum_{j\ge0}|\varepsilon_j|t^j\le(V-1)R(t)
 \qquad(t\ge0).
 \label{eq:pair-coeff-payment}
\end{equation}
Writing $(b_j)$ for the positive alphabet of $R$, every rectangular
reference determinant satisfies
\begin{equation}
 \Delta_{r,k}(R)\ge(b_1\cdots b_k)^r>0.
 \label{eq:pair-reference-margin}
\end{equation}
Let $\Theta=\sup_\ell\theta_\ell$, with $\Theta=0$ if there are no
nonreal pairs. If $(r+1)\Theta<\pi$, put
\[
 \kappa_r=\frac{\sin((r+1)\Theta)}{(r+1)\Theta},
 \qquad \kappa_r=1\text{ when }\Theta=0.
\]
Then
\begin{equation}
 \kappa_r^{rk}\le
 \frac{\Delta_{r,k}(E)}{\Delta_{r,k}(R)}
 \le\sec(\Theta)^{rk}.
 \label{eq:pair-sector-margin}
\end{equation}
\end{proposition}

\begin{proof}
For one pair, abbreviate its quadratic factor by $F$ and $(1+yz)^2$
by $G$. Direct algebra gives
\[
 s\cos\theta\le y<s,\qquad
 F-G=d\,z(z-1),\qquad 0<d/y^2\le\tan^2\theta.
\]
Thus the factors agree at zero and one. All products converge locally
uniformly. For $t\ge1$, $F(t)/G(t)$ increases from one to $s^2/y^2$.
Moreover
\[
 0\le2\log(s/y)\le-2\log\cos\theta\le\tan^2\theta.
\]
The summability assumption proves the finite quotient bound and its
limit by dominated convergence of the logarithms.

For finite products, expand the exact difference using
$F=G+d\,z(z-1)$. After taking coefficient absolute values, replace each
selected $z(z-1)$ by $z(z+1)$. At $t\ge0$ this gives the upper bound
\[
 R(t)\left\{\prod_\ell
 \left(1+\frac{d_\ell t(t+1)}{(1+y_\ell t)^2}\right)-1\right\}.
\]
Since $y\le2$ implies $y^2t(t+1)\le(1+yt)^2$, this is at most
$(V-1)R(t)$. Coefficient convergence and the same finite bound justify
passage to the infinite product.

The Toeplitz matrix of a finite positive-alphabet product is a product
of upper bidiagonal matrices $I+b_jS$. Their minors are nonnegative.
For the rectangle with row indices $0,\ldots,r-1$ and column indices
$k,\ldots,k+r-1$, Cauchy--Binet includes the positive term obtained by
shifting the whole block once in each of the first $k$ factors. Its
value is $(b_1\cdots b_k)^r$. Taking the infinite-product limit proves
\eqref{eq:pair-reference-margin}, including all negative-index zeros.

For completeness, rectangular Jacobi--Trudi identifies
$\Delta_{r,k}(E)$ with $s_{(r^k)}$ evaluated on its complete alphabet.
Classical Schur branching~\cite{EtingofRepresentationTheory20110110} into the real singleton and conjugate-pair
alphabets has nonnegative integer multiplicities. For a two-variable
partition $\nu=(u,v)$, $u\ge v\ge0$, the exact ratio is
\[
 \frac{s_\nu(se^{i\theta},se^{-i\theta})}{s_\nu(y,y)}
 =\left(\frac{s}{y}\right)^{u+v}
   \frac{\sin((u-v+1)\theta)}{(u-v+1)\sin\theta}.
\]
Each component appearing in the branching has $u\le r$: every
individual variable in $s_{(r^k)}$ has degree at most $r$, and the
nonnegative character multiplicities cannot cancel a highest-weight
monomial. Thus $u-v+1\le r+1$. Under the angular hypothesis, sinc
monotonicity, $\sin\theta\le\theta$, and $s/y\le\sec\Theta$ bound
each nonempty pair ratio below by $\kappa_r$ and above by
$\sec(\Theta)^{u+v}$. For the upper bound one also uses
$|\sin(j\theta)|\le j\sin\theta$, which follows from the finite
geometric sum. There are at most $rk$ nonempty pair components, and
their total degree is at most $rk$. Dividing the branching expansion
by its positive reference value expresses the full ratio as a
nonnegative weighted average of these products. This proves
\eqref{eq:pair-sector-margin}; local convergence gives the infinite
version.
\end{proof}

For the normalized theta function, the genus-zero product~\cite{Katkova0505174v1} is indexed
by upper-half-plane zeros $\rho=\sigma+i\gamma$ with multiplicity,
using $x=\alpha^{-1}$ and $\alpha=-(2\rho-1)^2$. The usual strip and
zero-counting inputs~\cite{Odlyzko1987AuthorCopy}, together with the inherited fact $\gamma>14$,
give $\Re x>0$, $|x|<1/784$, and
\[
 |\arg\alpha|=2\arctan\frac{|2\sigma-1|}{2\gamma},
 \qquad\sum\tan^2(\arg\alpha)<\infty.
\]
They verify the proposition's assumptions without RH or simplicity.
Both generating functions have nonnegative coefficients, so the
bounded positive-axis quotient also matches their maximum-modulus
growth up to a bounded factor. The reference nevertheless uses the
complete, implicit zero data; it is not an acquired numerical model.
No positive cosh-transform source or positive Green primitive for this
reference is asserted. The bounded-order estimate is classical sector
positivity, not a new theta positivity range. Requiring all pair Schur
blocks to be nonnegative at every degree would force every angle to
vanish and would restate RH for this product.

\begin{theorem}[Small complete scalar defects can have negative diagonal minors]
\label{thm:diagonal-pair-control}
Fix $B\in(1,2)$. For every integer $q\ge26$ there exist real entire
functions $E_q$ and $R_q$, with positive Taylor coefficients, such that:
\begin{enumerate}
 \item $E_q(0)=R_q(0)=1$ and $E_q(1)=R_q(1)=B$;
 \item both functions have genus zero, order $1/2$, and radial zero
 count asymptotic to $\sqrt T\log T/(8\pi)$;
 \item $R_q$ is an infinite product of positive linear factors, and
 $\Delta_{r,k}(R_q)>0$ for every $r,k\ge1$;
 \item writing $E_q-R_q=\sum_j\varepsilon_{q,j}z^j$, one has
 \begin{equation}
 \sup_{t\ge0}
 \frac{\sum_{j\ge0}|\varepsilon_{q,j}|t^j}{R_q(t)}
 \le\tan^2\!\left(\frac{9\pi}{4q+8}\right)=O(q^{-2}),
 \qquad \Delta_{q,q}(E_q)<0.
 \label{eq:diagonal-control-conclusion}
 \end{equation}
\end{enumerate}
The choices can additionally place every transformed zero in the
usual critical-strip geometry for all sufficiently large $q$.
\end{theorem}

\begin{proof}
Put $m=q-1$, $\theta=9\pi/(4q+8)$, and choose the unique $s>0$ solving
\begin{equation}
 (1+4s)^m(1+2s\cos\theta+s^2)=B.
 \label{eq:diagonal-endpoint}
\end{equation}
The left side increases strictly from one to infinity. Initially use
the finite alphabet
\[
 X=\{se^{i\theta},se^{-i\theta},\underbrace{4s,\ldots,4s}_{m}\},
 \qquad N=q+1,
 \qquad P=\prod_{x\in X}x>0.
\]
Finite elementary coefficient reversal, followed by transposition
and the one-column duality, gives
\begin{equation}
 \det[e_{q+j-i}(X)]_{i,j<q}=P^q h_q(X^{-1}).
 \label{eq:diagonal-reversal}
\end{equation}
Reversal is valid with all out-of-range coefficients zero; repeated
entries introduce no division. After removing the positive factor
$s^{-q}$, it suffices to sign
\[
 H_q=h_q(e^{i\theta},e^{-i\theta},1/4,\ldots,1/4).
\]
Let $Y$ have probability generating function
$((3/4)/(1-w/4))^m$. Its mean and variance are
$\mu=m/3$ and $W=4m/9$. Since the complete coefficients of the unit
pair equal $\sin((j+1)\theta)/\sin\theta$, coefficient extraction gives
\[
 H_q=\frac{(4/3)^m}{\sin\theta}
 \mathbb E\!\left[\sin((q+1-Y)\theta)\mathbf1_{\{Y\le q\}}\right].
\]
Now $(q+1-\mu)\theta=3\pi/2$. The inequality
$\cos x\ge1-x^2/2$ and Chebyshev's inequality imply
\begin{align*}
 \mathbb E[\sin((q+1-Y)\theta)\mathbf1_{\{Y\le q\}}]
 &\le-1+\frac{\theta^2W}{2}+\Pr(Y>q)\\
 &\le-1+\frac{\theta^2W}{2}+\frac{W}{(q-\mu)^2}
 <-1+\frac{13}{q}\le-\frac12.
\end{align*}
Indeed the two positive terms are less than $11.25/q$ and $1/q$,
respectively, using $\pi^2<10$. This makes
\eqref{eq:diagonal-reversal} strictly negative.

To complete the product, let
\[
 Q_J(z)=\prod_{j>J}(1+u_jz),\qquad
 u_j=\left(\frac{\log j}{4\pi j}\right)^2,\qquad J\ge3.
\]
Summability gives $Q_J\to1$ locally uniformly. For each fixed $q$,
choose $J$ sufficiently large that $Q_J(1)<B$, and solve
\eqref{eq:diagonal-endpoint} with $B/Q_J(1)$ in place of $B$, obtaining
$s_J\to s$. Set
\[
 E_q(z)=(1+4s_Jz)^m
 (1+2s_J\cos\theta\,z+s_J^2z^2)Q_J(z).
\]
This retains both endpoints. Its finite coefficient block converges
to the previously constructed block, so its $(q,q)$ determinant
remains strictly negative for all sufficiently large $J$.
This chooses an existent tail cutoff; no effective or certified
numerical cutoff is required for the assertion.

The tail zeros have moduli $(4\pi j/\log j)^2$. Inverting this sequence
gives the count $\sqrt T\log T/(8\pi)$. The convergence exponent is
$1/2$; also $u_j=O(j^{-2+\epsilon})$ for every $\epsilon>0$, so a
split-product bound gives order at most $1/(2-\epsilon)$. The count
gives the reverse order inequality. Thus the order is exactly $1/2$,
and summability gives genus zero. Finite factors change neither the
order nor the count asymptotic.

Replace the one conjugate pair by $(1+yz)^2$, with
$y=\sqrt{1+2s_J\cos\theta+s_J^2}-1$, and call the result $R_q$.
Its infinite positive alphabet gives all the strict reference margins.
The pair agrees at zero and one, and $y<\sqrt B-1<1$.
The exact difference is $d\,z(z-1)T(z)$, where
$d=s_J^2-y^2$ and $T=(1+4s_Jz)^mQ_J$ has positive coefficients.
Consequently
\[
 \frac{\sum_j|\varepsilon_{q,j}|t^j}{R_q(t)}
 \le\frac{dt(t+1)}{(1+yt)^2}
 \le\frac d{y^2}\le\tan^2\theta\qquad(t\ge0),
\]
which proves \eqref{eq:diagonal-control-conclusion}.

Finally $s\sim\log B/(4q)$. Choose $J(q)$ later, if necessary, so
that $s_J/s\to1$ while retaining the strict determinant sign. Then
\[
 \frac{\sin(\theta/2)}{\sqrt{s_J}}\to0,
 \qquad \frac{\cos(\theta/2)}{2\sqrt{s_J}}\to\infty.
\]
The reciprocal pair $\alpha=s_J^{-1}e^{\pm i\theta}$ can therefore
be written as $-(2\rho-1)^2$ with $0<\Re\rho<1$ for all sufficiently
large $q$. Positive real factors map to the critical line. This last
assertion concerns zero geometry only.
\end{proof}

The theorem concerns a family depending on $q$, not one entire function
proved to have infinitely many negative diagonal minors. Scaling both
functions by $c_0=b/B$ supplies the desired values $C(0)=c_0$ and
$C(1)=b$. No nonnegative smooth cosh-transform source, positive Green
primitive, full theta square-source identity, or ordinary Dirichlet
arithmetic is established for the controls. Thus the theorem excludes
a dimension-uniform positivity transfer using only the displayed data.
It does not exclude sharper dimension-dependent conditions, complete
signed compensation, or additional canonical source information.
The separate RGD-E04 r7 assessment accepted all twelve stated
obligations without mathematical correction. This is different-model
internal LLM assessment with shared premises, not specialist review.


\section{What remains unresolved}\label{sec:outlook}
The results support a comparison-methods paper, with a useful distinction
between an exact representation, a paid finite margin, and an estimate
that remains effective in growing dimension. They do not yet support an
all-order theta positivity theorem. The finite cases are already in known
positive territory; the conditional contour route leaves its scalar
constants and complementary region open; the fixed-index residue route
does not close the growing-index problem.

The pair construction supplies a particularly precise remaining question.
A determinant-sensitive estimate must use more than endpoints, leading
growth, positive coefficients and the scalar residual allowance retained
by the diagonal controls. Possibilities not excluded here include a
stronger dimension-dependent rate with explicit constants, signed
compensation between Schur branches, or a quantitative consequence of
the complete theta arithmetic source. No such estimate is supplied or
activated by this draft. Requiring every pair contribution to be
nonnegative at every degree would assume the zero reality one is trying
to prove.

The control distinctions are substantive. The positive cosh example
passes all rectangular signs while exhibiting the Gaussian scale
mismatch. A Gaussian-mixture example can retain an arbitrarily long finite
real-simple prefix of its base cosh function while adding nonreal zeros;
it is not asserted to match an actual theta zero prefix. The diagonal controls
retain the listed algebraic data but are not proved to have a nonnegative
smooth cosh-transform source, a positive Green primitive, or the canonical
theta square expansion. Other repository collision controls fail different
exact-source premises. A positive-time collision is not a counterexample
to every time-zero inequality. None of these examples proves a negative
canonical theta determinant.

The separate assessment of the pair construction is complete. Before
further mechanism design, this consolidated manuscript should be checked
as a whole. A later mathematical continuation should state which
information absent from the controls it will use and the quantitative
margin that information is expected to pay. Enlarging finite tables
without such a question would not close the gap identified here.

\appendix
\section{A conditional complete-contour margin}\label{sec:contour}
This appendix consolidates the conditional result of RGD Note08, whose
normalization and scope were checked in E04 r4. The canonical scalar
hypotheses below remain unproved here. No new positivity domain follows.

For $x>0$, let $\mu=xC'(x)/C(x)$ and
$v=\mathrm d\mu/\mathrm d\log x$. They are the mean and variance of
$c_jx^j/C(x)$. Positivity and the infinite coefficient support imply
$v>0$, $\mu(0+)=0$ and $\mu(x)\to\infty$; thus there is a unique
$x=x(k)$ with $\mu(x)=k$ for each integer $k\ge1$. Set
$\phi(\theta)=e^{-ik\theta}C(xe^{i\theta})/C(x)$.
Coefficient extraction and determinant integration give
\begin{equation}\label{eq:contour-int}
 D_{n,k}=\frac{C(x)^nx^{-nk}}{n!(2\pi)^n}
 \int_{[-\pi,\pi]^n}\prod_{i=1}^n\phi(\theta_i)
 \prod_{i<j}4\sin^2\frac{\theta_i-\theta_j}{2}\dd\boldsymbol\theta.
\end{equation}
This standard identity also appears in Katkova's Lemma~1
\cite{Katkova0505174v1}; the integral is real but not pointwise positive.

\begin{proposition}[Conditional joint-order allowance]\label{prop:contour}
Fix $0<c\le1$, $B>0$ and $0<\delta_0\le1$. At the saddle assume
\begin{align}
 \log\phi(\theta)&=-v\theta^2/2+R(\theta),&
 |R(\theta)|&\le Bv|\theta|^3\quad(|\theta|\le\delta_0),\label{eq:contour-H1}\\
 |\phi(\theta)|&\le e^{-cv\theta^2/2}&&(|\theta|\le\pi),\label{eq:contour-H2}
\end{align}
where the local logarithm is continuous and zero at the origin.
Put $A=8+2\log(2/c)$ and
\[
 G(n)=\max\left\{\frac{An^2}{c\delta_0^2},\,
 \frac{4An^4}{3c},\,16B^2(A/c)^3n^8\right\}.
\]
If $v\ge G(n)$, then
\begin{equation}\label{eq:contour-margin}
 D_{n,k}>\frac{C(x)^nx^{-nk}}{4n!(2\pi)^n}Z_n(v)>0,
 \qquad Z_n(v)=(2\pi)^{n/2}v^{-n^2/2}\prod_{j=1}^n j!.
\end{equation}
\end{proposition}
\begin{proof}
The Gaussian Vandermonde integral over $\R^n$ equals $Z_n(v)$.
Determinant integration and monic Hermite norms
$\sqrt{2\pi}\,j!v^{-j-1/2}$ prove this normalization, including the
$n!$ factor. Choose $\delta^2=An^2/(cv)$ and
$\tau=nBv\delta^3$. The assumptions give
$\delta\le\delta_0$, $\tau\le1/4$ and
$n(n-1)\delta^2/3\le1/4$. On $|\theta_i|\le\delta$, the real part
of the complete phase product is at least
$e^{-\tau}\cos\tau\,e^{-v\sum\theta_i^2/2}$.
The circular Vandermonde is at least the ordinary squared Vandermonde
times $S=(\sin\delta/\delta)^{n(n-1)}\ge e^{-1/4}$.
Here $\sin y/y\ge e^{-y^2/3}$ for $|y|\le1$, and all coincident
limits are continuous.

The Gaussian mass outside this cube is at most $TZ_n(v)$, where
$T=2^{n^2/2}e^{-v\delta^2/4}$. This follows by splitting the
Gaussian exponent in half on $\sum\theta_i^2\ge\delta^2$.
Using \eqref{eq:contour-H2}, the absolute contribution of the
\emph{entire} complementary contour region is at most $UZ_n(v)$, with
$U=(2/c)^{n^2/2}e^{-cv\delta^2/4}$, by the same argument at
variance parameter $cv/2$. This includes every mixed region; no
separate combinatorial exterior payment is omitted.
Our choice of $A$ gives $U=e^{-2n^2}$ and $T\le e^{-2n^2}$.
The real integral in \eqref{eq:contour-int}, divided by $Z_n(v)$, is
therefore at least
\[
 e^{-\tau}\cos\tau\,S(1-T)-U
 \ge e^{-1/2}\cos(1/4)(1-e^{-2})-e^{-2}
 >\frac{61}{224}>\frac14.
\]
This proves the claim.
\end{proof}

If canonical analysis were to establish the hypotheses for all $x\ge x_0$
with $\mu(x)/4\le v(x)\le\mu(x)$, then
$k\ge\lceil\max\{k_0,4G(n)\}\rceil$, for an integer
$k_0\ge\mu(x_0)$, would give an explicit sufficient domain.
Its $n^8$ scale is the cost of this allowance, not an optimal law.
The statement neither joins the complementary region nor constitutes
an induction at fixed $k$. The source literature already reports
stronger tail coverage; see the qualifications in the introduction.

\section{Supporting evidence and review status}\label{sec:provenance}
The companion file \texttt{CLAIM\_MAP.md} gives exact repository paths
for each assertion below. Notes, saved certificates and frozen assessment
requests remain the primary research record. This working manuscript
does not overwrite them or promote their status. The following describes
the record available when this draft was prepared.

\begin{center}\small
\begin{tabular}{>{\raggedright\arraybackslash}p{0.27\linewidth}>{\raggedright\arraybackslash}p{0.28\linewidth}>{\raggedright\arraybackslash}p{0.35\linewidth}}
\toprule
Material & Supporting record & Assessment status \\
\midrule
Green and finite comparison & Parent derivation; Notes01,03,05,06; saved E02 certificates & B01 accepted in E04 r1; two-case certificate accepted in r2; next-order exact record has author/same-model checks \\
Compact-reference stop and conditional contour & Notes07--09 & E04 r4 accepted within stated domains \\
Reciprocal comparison & Notes10--11 & E04 r5 accepted, twelve fixed obligations \\
Gaussian comparison & Notes12--13 & E04 r6 accepted, ten fixed obligations \\
Endpoint-pair reference and diagonal control & Notes14--15; Reviews11--12 & E04 r7 accepted, twelve fixed obligations, without mathematical correction \\
Manuscript integration & E01 r3 synthesis and integration check & Same-model editorial/proof-consistency checks; no new numerical run \\
\bottomrule
\end{tabular}
\end{center}

Acceptance here means an internal LLM assessment of the named contract,
with shared premises. The finite interval certificates depend on the
documented Arb/FLINT arithmetic and the analytical tail estimates; saved
record verification is not a fresh integration. The next-order result
uses existing moment enclosures and exact rational arithmetic. This
consolidation performs neither numerical acquisition nor a new zero
verification.

\paragraph{Source limits.}
The source ledger distinguishes inspected primary passages, inherited
published facts and unaudited reports. Platt--Trudgian's finite verification
is imported without replay. The reported fifteen-shift note was accessible
through indexed primary-page text when its direct page returned an error;
the supporting notebooks were not inspected. The original Schoenberg
article's metadata were checked, while the sector statement was checked
in Katkova and Khare. Etingof et al. supply the standard Schur-character
branching interpretation~\cite{EtingofRepresentationTheory20110110}.

Katkova's displayed Theorem~B uses the angle $\pi m/(m+1)$; a displayed
application uses $43\pi/44$ while labeling the conclusion PF$_{44}$.
The present work records this apparent integer mismatch without declaring
PF$_{44}$ false. It imports no numerical cutoff from that display and
proves its own symbolic angular condition. The r7 source check also notes
a mixed-sign factor display beside the conjugate quadratic on printed
p.~211 of Khare's notes; the self-contained pair proof here does not
use that display. Neither qualification is a correction to the present
formulas. A full literature and priority
audit remains outside this draft.

\paragraph{Reproducibility.}
The canonical editable files are in the investigation's
\texttt{manuscript/} directory. Build instructions and a small build
record accompany them. Generated working PDFs are stored in the
repository's ignored \texttt{output/pdf/} directory, not in manuscript
snapshot folders. Existing certificate scripts and small JSON records
remain in \texttt{numerics/}; external reading copies are governed by
their source receipts. The concise \texttt{DRAFT\_HISTORY.md} identifies
this first consolidated milestone without creating another snapshot.


\clearpage
\section*{Acknowledgments}
This manuscript was drafted for Edward Baker with substantial LLM
assistance in derivation, checking, exposition and document preparation.
The consolidation used gpt-6-astra at recorded effort ultra. Prior notes
and assessments identify their actual configurations, including
gpt-6.1-sol. Multiple model passes share source material and assumptions;
they do not constitute independent specialist verification. This is a working
manuscript for author review.

\bibliographystyle{plain}
\begingroup
\small
\bibliography{references}
\endgroup
\end{document}
